\section{引言}\label{sec:intro}

时间序列的结构变化既可能表现为均值、方差等低阶统计量的改变，也可能表现为局部轨迹几何、周期结构或状态空间形状的改变。滑动窗口把每个时间片段映射为点云，持续同调再从多尺度上提取连通分支与环，为后一类信号提供局部化表示\mcite{PereaHarer2015,Adams2017}；拓扑摘要已用于离线检测、在线监测与依赖序列中的结构异常识别\mcite{IslambekovYuvarajGel2020,ZhengMakXieXie2023,ChazalLevrardRoyer2024,NieYue2026}。为提高候选点密度，实践常取步长 $s<W$，相邻窗口因而共享原始观测，窗口特征序列随之形成局部依赖。

已有工作利用持续同调开展离线变点检测或在线监测；文献\cite{NieYue2026}进一步给出局部窗口映射的依赖继承结果并进行块校准。本文固定持续图向量表示与离线最大扫描，在两种观测预算下比较指定置换和依赖乘子规则的有限样本拒绝率，关注其逐条件偏差。

推断层面的困难来自随机化单位与科学零假设的层级错配。设第 $i$ 个窗口经固定映射得到特征 $Y_i$。无约束随机置换获得有限样本水平控制的常用理论依据，是窗口特征的联合分布对所用排列群保持不变：
\begin{equation}\label{eq:exchangeability}
 (Y_1,\ldots,Y_M)\overset{d}{=}(Y_{\pi(1)},\ldots,Y_{\pi(M)}),
 \qquad \forall\,\pi\in\mathfrak S_M .
\end{equation}
若窗口特征序列本身独立同分布，则式\eqref{eq:exchangeability}成立；原始观测的独立同分布性质不能在窗口重叠时自动传递为窗口特征的交换性。平稳性与相同边际分布也不足以一般地推出该联合对称性\mcite{RomanoTirlea2022,HemerikGoeman2018}。科学零假设定义在原始过程的“无变化”层面，置换参考分布所需的对称性则定义在派生特征的联合分布层面，二者并不自动一致。

本文评价一种固定实现：窗口点云经持续同调与固定向量化得到 800 维表示（BPI），或经均值—协方差得到低维表示（MOMENTS），再以标准化部分和的最大值作扫描检验，分别用无约束窗口级置换（M0）与依赖乘子自助法（M1）校准。本文在预设设计域内比较两种校准方法的有限样本表现，报告 28 个配置、112 个“配置—表示—方法”单元的拒绝率、Wilson 分类及三项预设对比；窗口均值反例与两类构造备择分别用于解释交换性风险和表示敏感性。本文不给出 BPI 扫描的通用有限样本定理，全部结论限定于所考察的设计网格与固定实现。
